\section*{Reading these results}
All forecasts concern monthly log realized variance, log(RV + 1e-8). The forecast evaluation spans June 1947 to December 2024, with an expanding unbalanced factor panel. Predictors begin in 1927. Training expands, validation uses 90 months, and models are selected each June.
The market portfolio is estimated both within the pooled factor panel and as a separate time series.
Pooled RMSE is the square root of the mean squared error across all series and months. It is therefore distinct from the average factor RMSE. R-squared is computed relative to the mean within each evaluation sample.
EW gives HARX, ENET, RF, XGB, NN1, NN2 and NN3 one-seventh each within each target-month, on the log scale. Each NN architecture counts once despite its ten internal members. Exponentiating this average is not the arithmetic mean of variance-level forecasts. No missing-member renormalization is allowed.
Validation scores describe selected hyperparameters on overlapping 90-month windows. Neural network tuning evaluates one network, while test forecasts average ten refitted networks. The comparison may also reflect refitting, window length, and changes in volatility conditions.
Feature summaries are descriptive linear surrogate loadings. Gross and net shares exclude the intercept. Exact collinearity can produce zero loadings and does not establish economic irrelevance. Predictor scales differ, and the equal-weighted forecast has no fitted surrogate model.
The sequential RMSE differences are arithmetic comparisons across model classes. They should not be interpreted as causal contributions. Elastic net adds shrinkage, while nonlinear learners also change regularization and estimation.
Portfolio Sharpe ratios and certainty-equivalent returns use excess returns. Cumulative wealth and maximum drawdowns use total returns, including the risk-free collateral return. Factor-panel portfolio returns are equal-weighted across the 153 factors within each month.
Managed exposure uses expanding volatility matching with at least 60 prior calibration observations. At each formation date, the normalizer matches unmanaged and inverse-variance-managed volatility over the available past observations. Exposure is capped at 1.5. Net returns deduct 10 basis points per unit change in factor exposure. These costs exclude trading within the underlying long-short factors.
Portfolio inference retains June 1952 to December 2024 for market strategies and January 1980 to December 2024 for the balanced factor portfolios. Each calendar uses shared month draws. The stationary bootstrap uses 10,000 draws, seed 20260918, and expected block length five, with three and ten as sensitivities. One-sided centered tests use the same prespecified families for Holm and BH adjustments. Spanning regressions use Newey-West lag one, with lags five and ten as sensitivities.
